Identifying the exact released versions affected by a known CVE is a core task for vulnerability management, downstream patch prioritization, and security database maintenance. Recent empirical evidence shows that existing affected-version identification tools remain brittle: tracing-based methods often propagate errors from vulnerability-inducing commit localization, while matching-based methods can miss semantically equivalent vulnerable code after refactoring, branch divergence, or non-deletion patches. This paper studies the problem setting of \emph{CVE affected-version identification}: given a CVE, a target project, its released versions, and the CVE fixing patch commits, recover the subset of released versions in which the vulnerable logic is still present.

We present \toolname, a repository-grounded framework that treats affected-version identification as evidence-constrained history reconstruction rather than as only commit-interval inference. The current method description is intentionally structured as a design placeholder: \toolname first filters and interprets the fixing commit family, then extracts root-cause-level vulnerability evidence, reconstructs boundary events, propagates vulnerability state through the Git DAG, and projects the result to released versions. We plan to evaluate \toolname on the benchmark introduced by ``Vulnerability-Affected Versions Identification: How Far Are We?'' and compare it against tracing-based and matching-based baselines under CVE-level exact accuracy, CVE-level no-miss ratio, and version-level precision/recall/F1. Full experimental results, ablations, and statistical analyses are \need{NEEDS EXPERIMENT}.
